\subsection*{Pertanyaan: Baca slide kuliah materi pencocokan string dengan algoritma Boyer-Moore di halaman 41. Tulis notasi algoritma untuk algoritma Boyer-Moore tersebut!}

\textbf{Notasi algoritma Boyer-Moore:}

\begin{lstlisting}[basicstyle=\ttfamily\scriptsize, frame=single, backgroundcolor=\color{white}, numbers=none, keywordstyle=\color{black}, commentstyle=\color{black}, stringstyle=\color{black}, breaklines=true]
Algoritma BoyerMoore(T, P)
Input  : teks T dengan panjang n, pola P dengan panjang m
Output : indeks awal P di dalam T, atau -1 jika tidak ditemukan

1. Bentuk tabel last occurrence:
   untuk setiap karakter c pada P,
   last[c] = indeks kemunculan terakhir c pada P

2. i = 0
3. Selama i <= n - m, lakukan:
      j = m - 1

      selama j >= 0 dan P[j] = T[i + j], lakukan:
          j = j - 1

      jika j < 0:
          kembalikan i

      karakter_salah = T[i + j]
      posisi_terakhir = last[karakter_salah]
      jika karakter_salah tidak ada pada P:
          posisi_terakhir = -1

      geser = maksimum(1, j - posisi_terakhir)
      i = i + geser

4. Kembalikan -1
\end{lstlisting}

\textbf{Tabel last occurrence untuk pola \texttt{acabaca}:}

\begin{table}[H]
    \centering
    \caption{Tabel Last Occurrence untuk Pola \texttt{acabaca}}
    \label{tab:praktikum10-last-occurrence-pretest}
    \begin{tabular}{cccc}
    \toprule
    \textbf{Karakter} & a & b & c \\ \midrule
    \textbf{Indeks terakhir} & 6 & 3 & 5 \\ \bottomrule
    \end{tabular}
\end{table}
